% Figure: the achromatic doublet. Left: a crown-glass positive element cemented
% to a flint-glass negative element brings the blue (F) and red (C) foci
% together, where a single lens would split them by chromatic aberration. Right:
% the residual focal shift against wavelength, crossing zero twice at the F and C
% design lines with a small secondary-spectrum bump between.
\begin{figure}[htbp]
\centering
\begin{tikzpicture}[>=Stealth, scale=1.0]
% ---- left panel: doublet ray schematic ----
\begin{scope}
% optical axis
\draw[enggray, dashed] (-3.4,0) -- (2.6,0);
% crown positive element (biconvex), drawn as a lens body
\fill[engblue!18] (-2.75,-1.15) .. controls (-2.35,0) .. (-2.75,1.15)
  .. controls (-2.35,0) .. (-2.75,-1.15) -- cycle;
\draw[engblue, thick] (-2.75,-1.15) .. controls (-2.35,0) .. (-2.75,1.15);
\draw[engblue, thick] (-2.75,-1.15) .. controls (-2.20,0) .. (-2.75,1.15);
% flint negative element cemented behind
\fill[engamber!18] (-2.10,-1.15) .. controls (-2.35,0) .. (-2.10,1.15)
  -- (-1.85,1.15) .. controls (-2.05,0) .. (-1.85,-1.15) -- cycle;
\draw[engamber, thick] (-2.10,-1.15) .. controls (-2.35,0) .. (-2.10,1.15);
\draw[engamber, thick] (-1.85,-1.15) .. controls (-2.05,0) .. (-1.85,1.15);
\node[engblue, font=\tiny, anchor=south] at (-2.75,1.18) {crown $(+)$};
\node[engamber, font=\tiny, anchor=south] at (-1.85,1.18) {flint $(-)$};
% incoming collimated rays (two heights)
\foreach \h in {0.55,-0.55} {
  \draw[engblack, thick] (-3.4,\h) -- (-2.55,\h);
}
% blue (F) rays converge to common focus at x=2.0
\draw[engblue, thick] (-1.95,0.55) -- (2.0,0);
\draw[engblue, thick] (-1.95,-0.55) -- (2.0,0);
% red (C) rays converge to the same common focus (achromat: F and C coincide)
\draw[engred, thick, dashed] (-1.95,0.55) -- (2.0,0);
\draw[engred, thick, dashed] (-1.95,-0.55) -- (2.0,0);
\filldraw[engblack] (2.0,0) circle [radius=1.3pt];
\node[font=\tiny, anchor=south west] at (2.0,0.05) {$F$ and $C$ focus};
\node[font=\scriptsize, anchor=north] at (-1.0,-1.35) {blue (F) and red (C) meet};
\end{scope}
% ---- right panel: focal shift vs wavelength ----
\begin{scope}[shift={(6.2,0)}]
\draw[->] (-0.3,-1.3) -- (-0.3,1.3) node[anchor=south, font=\tiny] {focal shift};
\draw[->] (-0.3,0) -- (3.2,0) node[anchor=west, font=\tiny] {$\lambda$};
% singlet curve: monotone (large chromatic aberration), for contrast
\draw[enggray, thick, domain=0.1:3.0, samples=40, variable=\x]
  plot ({\x}, {0.9 - 0.62*\x});
\node[enggray, font=\tiny, anchor=west] at (2.1,-0.75) {singlet};
% doublet curve: crosses zero twice (F and C), small hump between (secondary spectrum)
\draw[engblue, very thick, domain=0.1:3.0, samples=80, variable=\x]
  plot ({\x}, {0.30*(\x-0.7)*(\x-2.4)});
\node[engblue, font=\tiny, anchor=west] at (1.55,0.55) {doublet};
% mark F and C zero crossings
\filldraw[engred] (0.7,0) circle [radius=1.1pt];
\filldraw[engred] (2.4,0) circle [radius=1.1pt];
\node[engred, font=\tiny, anchor=north] at (0.7,-0.06) {F};
\node[engred, font=\tiny, anchor=north] at (2.4,-0.06) {C};
\end{scope}
\end{tikzpicture}
\caption[The achromatic doublet]{A converging crown-glass element cemented to a
diverging flint-glass element brings the blue (Fraunhofer F) and red (C) foci
together, cancelling first-order chromatic aberration where a single lens (grey,
right panel) spreads the focus steadily with wavelength. The doublet's residual
focal shift crosses zero at the two design lines and leaves only a small
secondary spectrum between them, the bump an apochromat removes by adding a third
glass. The crown carries the net positive power and the flint carries the
dispersion that cancels the crown's, so the pair is solved from the two element
powers and the two Abbe numbers together.}
\label{fig:vol04ch11:achromat-doublet}
\end{figure}
